\documentclass[11pt]{article}
\usepackage[margin=1in]{geometry}
\usepackage{amsmath,amssymb,amsthm,hyperref}
\newtheorem{proposition}{Proposition}
\newcommand{\E}{\mathbb E}
\newcommand{\F}{\mathbb F}
\newcommand{\Prob}{\mathbb P}
\title{Joint target existence in independently hashed elliptic-curve slots}
\author{Project research note}
\date{20 November 2026}
\begin{document}
\maketitle

\paragraph{Result.}
A single determinant represents a target hit on regular inputs without any intermediate curve additions.
Elliptic-curve interpolation computes all determinant moments through order
$\lfloor 2^{b/m^2}\rfloor$ over the whole Cartesian product from the full
hash seeds within the desired time budget. On prime-order curves, even
the combined zero tests of all these moments have fresh error of order
at least $1/\log b$ averaged over eligible targets. The reduction requires
$O(1/(\log b)^2)$ for every
prescribed target. The end-to-end exponent and its leading constant stay at
the previous conditional bound.

\section{The determinant and the fresh-input error}

Write $G=E(\F_p)$, $N=hr=|G|$, $p>3$, with the project's supplied coprime
cofactor and odd prime subgroup of order $r$. Let $O$ be the identity and
$T$ a nonzero target in that subgroup. Put $b=\lceil\log_2 p\rceil$,
$M=2^b$, $L=\log_2 b$, $g=\log_2 L$. The slot $F_i$ contains all affine
points whose canonical abscissa $x$ satisfies
$(\alpha_i x+\beta_i\bmod M)<\tau_i$. The multipliers are uniform odd,
the offsets are uniform, and all slots are independent.
Use the balanced dyadic widths of \texttt{fresh\_windows.tex},
$\rho_i=\tau_i/M$, $\nu_i=N\rho_i$, and
$1/(2L)<\lambda=N^{-1}\prod_i\nu_i\le1/L$.
The proposed arity is $m\asymp L^{1/4}\sqrt g$.
Write $H_T$ for the ordinary signed, ordered count of tuples of sum $T$.
The targetwise moments in \texttt{../mixed\_coefficient/mixed.tex} give
$\Prob(H_T=1)=\lambda+O(\lambda^2)$ and
$\Prob(H_T\ge2)=O(\lambda^2)$ uniformly at these widths.

Let $d=m+1$ and use the monomial basis of $L(dO)$ consisting of $x^j$
with $0\le 2j\le d$ and $yx^j$ with $0\le 2j+3\le d$, the constant
first. Let $v(P)$ be its evaluation column at $P$ and set
\[
 \Delta_{\mathbf P}(T)=\det[v(P_1),\ldots,v(P_m),v(-T)] .
\]
For distinct affine points in these $d$ positions, $\Delta=0$ if and only if
$\sum_iP_i=T$. In fact a kernel of the evaluation map is a nonzero section
vanishing at those $d$ points. Its divisor must be their sum minus $dO$,
and on an elliptic curve a degree-zero divisor is principal exactly when
its group sum is $O$. Conversely, a principal divisor gives the required
section, unique up to a scalar. This is the usual elliptic-code evaluation
criterion. See \href{https://arxiv.org/html/cs/0507026}{Cheng,
\emph{Hard problems of algebraic geometry codes}, Sections 2 and 3}.

Call the seeds regular for $T$ if slots have pairwise disjoint abscissa sets
and avoid $x(T)$. Each abscissa is kept with probability $\rho_i$.
Consequently failure has probability at most
\[
 N\sum_{i<j}\rho_i\rho_j+\sum_i\rho_i=o(\lambda^2)
\]
at the stated widths. Here $\log_2\nu_i=b/m+O(1+(\log L)/m)$ and $m$ grows.
One may explicitly test regularity during slot construction and return zero
on failure. Thus existence of a singular choice gives a fresh parity
predictor of error $O(\lambda^2)$ \emph{if one can compute it}.
Testing a determinant for a prescribed tuple costs polynomial time, but
the expected number of candidate tuples is
$\prod_i\E|F_i|=(1-1/N)^m N\lambda$.

\section{A joint statistic that fits the time budget}

Assign optional field weights $w_{i,P}$ to points of each slot. For $j\ge1$
define
\[
 Z_j(T)=\sum_{\mathbf P\in\prod_i F_i}
       \Delta_{\mathbf P}(T)^j\prod_i w_{i,P_i}.
\]
With $v(-T)$ as fixed column, subtract it from each other column and omit
the constant coordinate to obtain $m$-vectors $q_T(P)$. This changes at
most the sign of the determinant. Put
$G_i(T)=\sum_{P\in F_i}w_{i,P}q_T(P)q_T(P)^{\mathsf t}$.
Then Cauchy--Binet gives
\[
 Z_2(T)=[z_1\cdots z_m]\det\left(\sum_i z_i G_i(T)\right)
 =\sum_{S\subseteq[m]}(-1)^{m-|S|}\det\left(\sum_{i\in S}G_i(T)\right).
\]
Indeed the squarefree term selects precisely one column from each slot.
The subset formula follows by inclusion-exclusion since the determinant is
homogeneous of degree $m$. Empty slots give zero. The computation takes
$O(m^2\sum_i |F_i|+2^m m^3)$ field operations using subset sums or Gray
order, in addition to the cost of constructing the slots and their weights.
Enumerating $\tau_i$ inverse-hash residues, discarding those outside the
field and lifting the remainder uses expected $\operatorname{poly}(b,m)
\sum_i\tau_i$ bit operations. Unit or independent uniform point weights
fit this accounting. Thus the base-two logarithm of the expected bit cost
is at most $b/m+O(m+\log b)$, including setup. At the proposed arity
this meets the scale required for the \emph{evaluation} subroutine.

One can also preprocess all seeds independent of $T$. Set
$A_i=\sum_{P\in F_i} w_{i,P}v(P)v(P)^{\mathsf t}$, of size $d$.
The rank-one determinant identity gives the formula valid even for singular
matrices
\[
 Z_2(T)=v(-T)^{\mathsf t} C v(-T),\quad
 C=\sum_{S\subseteq[m]}(-1)^{m-|S|}
             \operatorname{adj}\left(\sum_{i\in S}A_i\right).
\]
To verify it, extract $z_0 z_1\cdots z_m$ from
$\det(z_0 v(-T)v(-T)^{\mathsf t}+\sum_i z_i A_i)$.
This describes the entire target dependence of the second moment with
$d^2$ field entries. With unit weights, summing over $P$ and $-P$
cancels every mixed even/odd coordinate in $A_i$. The adjugates are also
block diagonal (by the cofactor formula). On $y^2=f(x)=x^3+ax+c$,
$Z_2$ therefore becomes a polynomial in $x(T)$ of degree at most $d$.
A point of order two contributes zero to each mixed coordinate as well.

\section{Computing many higher moments jointly}

\begin{proposition}
Let $k\ge1$, $n=kd$, with $n+1$ at most half the affine curve population.
Given unit weights, or field weights supplied at unit field-operation cost
per slot point, all $Z_1(T),\ldots,Z_k(T)$ can be evaluated from the
seeds in expected bit cost at most
\[
 \operatorname{poly}(b,m,\log(k+1))
 \left(n\sum_i\tau_i+n^3+mn^2+kn^m\right).
\]
In particular $k=\lfloor 2^{b/m^2}\rfloor$ gives log-cost at most
$(1+O(1/m))b/m+O(m\log(m+1)+\log b)$.
\end{proposition}
\begin{proof}
Construct $n+1$ distinct affine curve points $B$ by sampling field
abscissas and lifting. Hasse's estimate, at most two points per abscissa,
and the stated occupancy bound give expected $\operatorname{poly}(b)$
bit operations per new point (one may sample uniformly from signed affine
points by rejection on an abscissa and a sign). Their evaluation matrix
for the $n$ monomials of $L(nO)$ has full column rank: a nonzero section
cannot vanish on $n+1$ distinct points. Select $n$ independent rows by
elimination and denote their points $B_1,\ldots,B_n$.

For each slot compute its weighted sums of all $n$ monomials by tabulating
powers of $x$ and multiplying by $y$ where required. Using the common
evaluation matrix, solve for weights $h_{i,a}$ such that
\[
 \sum_{P\in F_i}w_{i,P}\,\phi(P)=\sum_{a=1}^n h_{i,a}\phi(B_a)
 \qquad\text{for every }\phi\in L(nO).
\]
For any $j\le k$, holding the other arguments fixed,
$\Delta_{\mathbf P}(T)^j$ belongs to $L(jdO)\subseteq L(nO)$ in the
remaining slot. Apply the identity one slot at a time. This proves
\[
 Z_j(T)=\sum_{a_1,\ldots,a_m=1}^n
  \left(\prod_i h_{i,a_i}\right)
  \Delta_{B_{a_1},\ldots,B_{a_m}}(T)^j .
\]
Compute the determinant once per node tuple and accumulate its first
$k$ powers by successive multiplications. Slot enumeration, monomial
sums, the common factorization with $m$ right-hand sides, and this final
sum give the claimed costs. With the displayed choice of $k$,
$\log_2 n\le b/m^2+\log_2 d$ and
$\log_2(\sum_i\tau_i)\le b/m+O(\log b)$.
\end{proof}

This is a joint computation on a smaller, public grid. It uses field
weights on its nodes, so testing those nodes individually for singularity
does not decide whether any original tuple is singular. The final sum
preserves each specified power moment.

\section{How accurately can zeros of these moments predict a hit?}

The following statement allows the coefficients of the tests to be computed
from \emph{all} the hash seeds, at arbitrary cost, before the target is chosen.

\begin{proposition}
For fixed seeds and optional random tape, suppose the prediction is a Boolean
combination of zero tests of functions on $E$ with their only pole at $O$
and total pole-order bounds summing to $D$. The functions and Boolean rule
are fixed before varying $T$. In particular one can use $Z_j-c$
at a pole-order charge $jd$, with weights and comparison values depending
arbitrarily on the seeds, independently of $T$. For uniform nonzero $T$
in the subgroup of order $r$, independent of the seeds, error against
existence or parity is at least
\[
 \lambda-O(\lambda^2)-\frac{D+2\lambda}{r-1}
\]
in the tuned regime. The displayed bound applies for a uniform upper bound
on $D$, or with $\E D$ in its place.
\end{proposition}

\begin{proof}
An identically zero function has a constant zero-test bit. Any other
function of pole order at most $k$ has at most $k$ zeros. For the fixed
seeds, the prediction agrees with a constant off at most $D$ subgroup
points. The bound on $Z_j$ follows by expanding the determinant along the
target column. The map $T\mapsto -T$ preserves $O$, and products of $j$
sections of $L(dO)$ belong to $L(jdO)$.

For completeness let $u$ be the fraction of eligible targets positive
for the true bit at these seeds. The targetwise moments give
$\E u=\lambda+O(\lambda^2)$ for either truth definition.
Write $A_0=\prod_i(1+2/\nu_i)$ and $z=|G[2]|(2^m-2)/N$.
We claim
\[
 \E u^2\le A_0\lambda^2+z\lambda+2\lambda/(r-1).
\]
Bound both positivity bits by their ordinary counts and take two independent
uniform targets $T,S$. For a pair of tuples, classify the indices where
their abscissas coincide. If this set $J$ is proper, there are at most
$N^{m-1}2^{|J|}N^{m-|J|-1}$ pairs for the two prescribed sums. The
two-abscissa hash law bounds the retention probability by
$\prod_{i\in J}\rho_i\prod_{i\notin J}\rho_i^2$. Summing gives the first
term. When all abscissas agree, a nonempty proper sign-flip set $I$
forces $2\sum_{i\in I}P_i=T-S$. This permits at most $|G[2]|N^{m-2}$
first tuples, giving the second term. Empty and full flips require
$S=T$ and $S=-T$, respectively. Each contributes at most $\lambda$
when it occurs. Averaging gives the claim. The risk of any constant bit
is at least $\min(u,1-u)\ge u-u^2$. Changing its predictions on $D$
targets can reduce risk by at most $D/(r-1)$. Now $A_0=1+o(1)$ and
$z=o(\lambda)$, proving the result.
\end{proof}

For $Z_2=0$, $D=2d$. On the prime-order class, and more generally for
$d/r=o(\lambda)$, the error is at least $(1-o(1))\lambda$. Thus there
is a prescribed target at which the required fresh performance fails.
For a family of moment zero tests to work uniformly, a necessary condition
in this model is $\E D\ge (r-1)(\lambda-O(\lambda^2))-2\lambda$.
A large degree can have a small circuit. The argument does not imply a
corresponding runtime lower bound and does not cover arbitrary ways of
combining a particular target and the full seeds.

In particular, testing each of the computed $k=\lfloor 2^{b/m^2}\rfloor$
moments for zero gives $D\le d k(k+1)/2$. On prime-order curves this is
$o(r\lambda^2)$. Even an arbitrary Boolean decision based on all these
zero bits, with the rule chosen from the seeds before varying $T$, has
target-averaged error at least $(1-o(1))\lambda$. Decoding the entire
vector of field values is a different question and remains open.
The bound also reaches nonlinear tests on those values. For example,
take the determinant of the $s$-by-$s$ matrix $(Z_{i+j})_{0\le i,j<s}$,
with $Z_0$ the total weight and $2s-2\le k$. Its pole-order bound is
$d\,s(s-1)$ since every term of the determinant has the same sum of
moment indices. Testing all such leading minors costs $O(dk^3)$ in the
degree bound. This too is $o(r\lambda^2)$ on the prime-order class.

There is a separate per-target test for independent uniform weights.
On any fixed input admitting a nonsingular tuple, regard $Z_2$ as a
linear form in the weights of one slot, after conditioning on all other
weights. At least one coefficient is a nonzero polynomial of degree
$m-1$ in those other weights: the nonsingular tuple supplies a distinct
nonzero monomial. The Schwartz--Zippel bound shows that the form is
nonzero except with probability at most $(m-1)/p$. Thus the scalar alone
is within $(m-1)/p$ in total variation of the uniform field element.
Typically this holds on both positive and negative fresh inputs.
Indeed, distinct abscissas have nonpositive indicator covariance;
with at most two points per abscissa,
$\operatorname{Var}|F_i|\le 2\E|F_i|$. All slots have more than $2d$
points except with probability $O(\sum_i1/\nu_i)=o(\lambda^2)$.
Together with regularity this allows any choice in the first $m-1$ slots
and multiple choices in the last, at most one of which can sum to $T$.
Consequently a decoder seeing only the randomized scalar has error at least
$\lambda-O(\lambda^2)$ at \emph{each} fixed target. This last bound does
not include additional information about the realized weights or seeds.

\section{The high moment that would suffice}

On regular inputs, take all weights equal to one. Fermat's identity gives
\[
 H_T\bmod p=\left(\prod_i |F_i|-Z_{p-1}(T)\right)\bmod p .
\]
Outputting whether this residue is nonzero gives error $O(\lambda^2)$
for ordinary parity, since it is exact when $H_T\le1$. The same error
bound holds for predicting ordinary existence. The low cost of the
Cauchy--Binet calculation for $Z_2$ has not been established for this
moment. Raising an already summed statistic to a power introduces cross
terms and does not compute the sum of powers. Raising \emph{each}
determinant to $p-1$ is inexpensive on a single tuple, but enumerating
all tuples costs order $N\lambda$ in expectation before field arithmetic.

There is a specific reason the interpolation algorithm cannot simply
use a short list of nodes at this order. Suppose a slot sum with nonzero
weights on $F$ is to be represented by weighted point evaluations on
an affine node set $B$, exactly on all of $L(nO)$. If
$n>|F\cup B|$, then the evaluations at $F\cup B$ are surjective:
Riemann--Roch gives kernel dimension $n-|F\cup B|$ for that evaluation
map. Hence any point of $F$ omitted by $B$ is distinguished by a section
equal to one there and zero at every other point of the union. Exactness
forces $F\subseteq B$. At $n=d(p-1)$ the inequality holds even for the
entire affine population at the tuned arities. Thus preserving this
whole function space recovers the original slot points. An evaluator
tailored just to the determinant power might use less information.

\paragraph{Sampling a proposed divisor.}
Another encoding chooses the section with the target among its zeros.
For distinct abscissas $x_i$ avoiding $x_T=x(T)$, write the section as
$U(x)+yV(x)$. Having chosen one sign above each $x_i$, the group sum
can equal $T$ exactly when there is a normalized section with
\[
 U(x_T)=y_TV(x_T),\quad V(x_T)\ne0,\qquad
 U(X)^2-f(X)V(X)^2=c_d(X-x_T)\prod_i(X-x_i).
\]
For $d=2n$ take $U$ monic of degree $n$, $\deg V\le n-2$, $c_d=1$.
For $d=2n+1$ take $V$ monic of degree $n-1$, $\deg U\le n$, $c_d=-1$.
The squarefree right-hand side precludes any common root of $U,V$ at
these abscissas. The section selects the signed ordinates $-U/V$ and has
exactly the required zero divisor. Conversely the divisor supplies such a
section by normalization. This equivalence includes an abscissa of a
point of order two when the roots on the right are simple.

Testing only $s\le\lfloor d/2\rfloor-2$ distinct field nodes in addition
to the signed target is always satisfiable if the polynomials may be
chosen after seeing the nodes. Take nodes different from $x_T$ with
$f(z)\ne0$. At each node, $u^2-f(z)v^2$ represents every field element
by a split product or the surjective quadratic extension norm. Prescribe
one such local pair, then prescribe $V(x_T)=1$, $U(x_T)=y_T$.
Each polynomial has at least $n-1$ free coefficients after fixing its
leading term where necessary. Interpolation at $s+1\le n-1$ nodes fits
these data. This rejects the proposed sampled existential equations even
on negative inputs. A challenge sampled after committing to a whole
section is a different computation; the cost of choosing that section
would still need to be included.

\paragraph{Batching the existing chord arithmetic.}
Consider directly replacing the $e=2m-2+t$ alternating-chord equations
by $s$ uniform independent field-linear combinations, keeping the
canonical word ranges, the hash tests and the unequal-abscissa guards.
Here $t=\lfloor(m-1)/2\rfloor$ is the number of free internal even
ordinates. Set every slope to one. Pick legal abscissas, using disjoint
nonempty slots, and canonical internal abscissas satisfying the guards.
These choices exist on the typical regular seeds above. In the retained
ordinate vector $Y=(Y_2,Y_4,\ldots)$ the remaining residuals are
$R(Y)=R_0+JY$. In the even step $j<m$, the row $D_j$ contains $2Y_j$
and none of the other retained ordinates. This follows directly from the
alternating chord equations in the project's compressed-relation manuscript.
Thus $J$ contains a $2I_t$ minor and has column rank $t$.

For a uniform $s$-by-$e$ field matrix $C$, $CJ$ is a uniform
$s$-by-$t$ matrix. With probability $\prod_{k=0}^{s-1}(1-p^{k-t})$ it
has full row rank for $s\le t$. Then $CJY=-CR_0$ has a solution for
every $R_0$, even on negative seeds. The proposed $s=O(m/g)$ eventually
satisfies $s\le t$, so the compressed equations accept almost all fresh
negative inputs. This applies to batching the residuals in the original
encoding, which derives intermediate curve membership from the complete
equations. It does not cover additional independent constraints or
tests performed after a complete witness is fixed.

\paragraph{Enumerating completions.}
Suppose a procedure sees all seeds but the last offset (it may also see
that slot's multiplier), prepares an explicit list of $K$ possible
completion points, and returns a positive only on a verified listed hit.
Each listed affine point has conditional acceptance probability $\rho_i$.
At $O(\lambda^2)$ error one must detect positives of mass at least
$\lambda-O(\lambda^2)$, so
$\E K\ge(1-O(\lambda))\lambda/\rho_i$. This is sharp to relative
$O(\lambda)$: enumerating the other-slot tuples and listing their
nonidentity required completions has expected length $\E H_T/\rho_i$.
At the proposed widths the list has logarithmic size
$(1-1/m)b+O(\log b)$. The statement concerns explicit completion lists
constructed without the last offset. A computation with all offsets jointly
is still possible in principle.

\paragraph{End-to-end accounting.}
The outer reduction charges polynomial factors times $B_0 C_m+mB_0^2$,
where $B_0=2\sum_i\tau_i$ and $C_m$ includes verified row selection.
Under the fresh-error reduction in \texttt{fresh\_windows.tex}, an evaluator
with $\log_2 D_m=O(mb/(\sqrt L\,g))$, including preprocessing, would give
$\log_2\E T_{\rm solve}=O(b/(L^{1/4}\sqrt g))$ at the proposed arity.
The fast moment misses the error budget. The exact high-moment predicate
meets the error budget but has no evaluation at that cost.
The complete bound remains
$\log_2\E T_{\rm solve}\le(26.2933775156+o(1))b/L^{1/4}$
under the finite-packet assumption A.1 and the inherited contraction and
reduction estimates of \texttt{../short\_hash\_count/count.tex}.

\paragraph{Reproducibility.}
Run \texttt{python -B research/joint\_parity/check\_target\_incidence.py}
from the project root. The script checks all 12,096 signed distinct-abscissa
configurations of degrees three through six on an order-19 curve (630 singular).
It also enumerates 294,912 seed-pair/target cases, checks both weighted
and unweighted Gram formulas, checks the Fermat residue on regular inputs,
and exhibits regular inputs with the same target, slot sizes and second
moment but opposite existence bits. A separate order-103 example passes a
sampled norm equality although all 32 sign choices miss the stated target.
On that same curve, the weighted interpolation calculation through power
five agrees between 8,000 public node tuples and 16,896 original tuples.
The same negative example also tests random batches of the project's
actual chord equations. The error bounds above follow from the
asymptotic arguments.
\end{document}
